% GENERATED FILE. Edit canonical lessons, then run npm run generate:books.
% canonical-source-hash: fnv1a64:20f2f96bc5b1d657
% canonical-lessons: AR-C286-an-adhhaba, AR-C286-yajibu-an, AR-C286-bi-sabab, AR-C286-sawfa, AR-R286-advice-and-reasons

\chapter{Must, Should, Will, Because}
\label{ch:ar-must-should-will-because}
\hlchaptermodality{286}

\begin{chapteropening}
\textbf{By the end of this chapter:} \emph{I can say what I want, must or ought to do, give reasons for my opinions and plans, and say what I will and will not do.}

Complete AR-R286-advice-and-reasons: a plan and a piece of advice with the subjunctive, and the reasons for both.
\end{chapteropening}

\section[\texorpdfstring{\ar{أن} \ar{أذهب} (an adhhaba)}{أن أذهب (an adhhaba)}]{\ar{أن} \ar{أذهب} (an adhhaba) — to go: the verb after an}

\label{lesson:AR-C286-an-adhhaba}



\begin{quote}
Say \textquotedblleft{}I think that the exam is hard\textquotedblright{}. (\emph{Aʿtaqidu anna al-imtiḥāna ṣaʿb}.) Then \textquotedblleft{}in my opinion\textquotedblright{}. (\emph{Fī raʾyī}.)

You have said \emph{anwī an azūra jaddatī}, \textquotedblleft{}I intend to visit my grandmother\textquotedblright{}. Listen to the end of \emph{azūra}. It is not the \emph{azūru} of \textquotedblleft{}I visit\textquotedblright{}, and this lesson says why.
\end{quote}

\begin{grammarlens}[title={Grammar Lens: -u becomes -a}]
After \textbf{\ar{أن}} — \emph{an}, \textquotedblleft{}to\textquotedblright{}, \textquotedblleft{}that\textquotedblright{} before a verb — the present changes its last vowel. \emph{Adhhabu}, I go, becomes \emph{an adhhaba}, to go:

\begin{quote}
\textbf{\ar{أريد} \ar{أن} \ar{أذهب} \ar{إلى} \ar{السوق}.} — \emph{urīdu an adhhaba ilā as-sūq} — I want to go to the market.
\end{quote}

This shape of the present is the \textbf{subjunctive}: the verb for a thing wanted, planned or needed, not yet a fact. \emph{Urīdu}, I want, is the present of \emph{arāda}.

Where the present ends in \emph{-ūna} or \emph{-īna}, the \emph{-na} falls away: \emph{yadhhabūna}, they go, becomes \emph{an yadhhabū}; \emph{tadhhabīna}, you go (to a woman), becomes \emph{an tadhhabī}.
\end{grammarlens}

\begin{grammarlens}[title={Grammar Lens: an and anna}]
\begin{quote}
\textbf{\ar{قررنا} \ar{أن} \ar{نسافر} \ar{بالقطار}.} — \emph{qarrarnā an nusāfira bi-l-qiṭār} — We decided to travel by train.
\end{quote}

\emph{An} takes a verb: what someone wants or decides to do. \emph{Anna}, with the doubled \emph{n}, takes a sentence: what someone thinks is so. \emph{Aʿtaqidu anna al-qiṭāra sarīʿ}; \emph{urīdu an arkaba al-qiṭār}.
\end{grammarlens}

\subsection*{Your turn}
Say these aloud:

\begin{itemize}
\raggedright
  \item \emph{adhhabu} … \emph{an adhhaba}
  \item \textquotedblleft{}I want to go to the market\textquotedblright{} — \emph{urīdu an adhhaba ilā as-sūq}
  \item \textquotedblleft{}We decided to travel by train\textquotedblright{} — \emph{qarrarnā an nusāfira bi-l-qiṭār}
  \item \textquotedblleft{}They want to go\textquotedblright{} — \emph{yurīdūna an yadhhabū}
  \item \emph{anwī an azūra jaddatī}, and hear the \emph{-a}
\end{itemize}

\begin{tcolorbox}[breakable,colback=teal!4,colframe=teal!35!black,title={Before you move on}]
After \emph{an}, what happens to \emph{-u}? (\textbf{It becomes \emph{-a}}: \emph{an adhhaba}.) And to \emph{-ūna}? (\textbf{The \emph{-na} falls away}: \emph{an yadhhabū}.) Which particle takes a verb, and which a sentence? (\emph{\textbf{An}} a verb, \emph{\textbf{anna}} a sentence.)
\end{tcolorbox}
\section[\texorpdfstring{\ar{يجب} \ar{أن} (yajibu an)}{يجب أن (yajibu an)}]{\ar{يجب} \ar{أن} (yajibu an) — must, should}

\label{lesson:AR-C286-yajibu-an}



\begin{quote}
Say \textquotedblleft{}I want to go\textquotedblright{}. (\emph{Urīdu an adhhaba}.) Then the Arabic for homework, a duty. (\emph{Wājib}.)

Wanting is one reason to act. Having to is another, and a piece of advice is a third.
\end{quote}

\begin{grammarlens}[title={Grammar Lens: yajibu an, it is necessary}]
\textbf{\ar{يجب}} — \emph{yajibu} — \textquotedblleft{}it is necessary\textquotedblright{}. It does not change for the person. The person is in the verb after \emph{an}, in the subjunctive:

\begin{quote}
\textbf{\ar{يجب} \ar{أن} \ar{أذهب} \ar{الآن}.} — \emph{yajibu an adhhaba al-ān} — I must go now.
\end{quote}

So \emph{yajibu an tadhhaba}, you must go; \emph{yajibu an nadhhaba}, we must go. Put \emph{kāna} in front for \textquotedblleft{}had to\textquotedblright{}: \emph{kāna yajibu an adhhaba}, I had to go.

\emph{Yajibu} comes from the root of \emph{wājib}, \textquotedblleft{}a duty\textquotedblright{}. The root's \emph{w} drops in the present, as it does in \emph{yaṣilu}, he arrives, from \emph{waṣala}.
\end{grammarlens}

\begin{grammarlens}[title={Grammar Lens: yanbaghī an, you ought to}]
\textbf{\ar{ينبغي}} — \emph{yanbaghī} — \textquotedblleft{}it is right, it is fitting\textquotedblright{}. It is softer than \emph{yajibu}: advice, not an order.

\begin{quote}
\textbf{\ar{ينبغي} \ar{أن} \ar{تزور} \ar{الطبيب}.} — \emph{yanbaghī an tazūra aṭ-ṭabīb} — You ought to visit the doctor.
\end{quote}

\emph{Yanbaghī an tastarīḥa}, you should rest. The verb after it behaves exactly as after \emph{yajibu}.
\end{grammarlens}

\subsection*{Your turn}
Say these aloud:

\begin{itemize}
\raggedright
  \item \textquotedblleft{}I must go now\textquotedblright{} — \emph{yajibu an adhhaba al-ān}
  \item \textquotedblleft{}We must go\textquotedblright{} — \emph{yajibu an nadhhaba}
  \item \textquotedblleft{}I had to go\textquotedblright{} — \emph{kāna yajibu an adhhaba}
  \item \textquotedblleft{}You ought to visit the doctor\textquotedblright{} — \emph{yanbaghī an tazūra aṭ-ṭabīb}
  \item advice with a reason — \emph{fī raʾyī yanbaghī an tastarīḥa, li-annaka marīḍ}
\end{itemize}

\begin{tcolorbox}[breakable,colback=teal!4,colframe=teal!35!black,title={Before you move on}]
Does \emph{yajibu} change for the person? (\textbf{No}: the verb after \emph{an} does.) The softer one, for advice? (\emph{\textbf{Yanbaghī.}}) \textquotedblleft{}I had to go\textquotedblright{}? (\emph{\textbf{Kāna yajibu an adhhaba.}}) Which noun shares the root of \emph{yajibu}? (\emph{\textbf{Wājib}}, a duty.)
\end{tcolorbox}
\section[\texorpdfstring{\ar{بسبب} (bi-sabab)}{بسبب (bi-sabab)}]{\ar{بسبب} (bi-sabab) — because of}

\label{lesson:AR-C286-bi-sabab}



\begin{quote}
Say \textquotedblleft{}I must go now\textquotedblright{}. (\emph{Yajibu an adhhaba al-ān}.) Then ask \textquotedblleft{}why?\textquotedblright{}. (\emph{Limādhā?})

The answer can be a whole sentence, or a single thing. Arabic uses a different word for each.
\end{quote}

\begin{grammarlens}[title={Grammar Lens: a thing, or a sentence}]
\textbf{\ar{بسبب}} — \emph{bi-sabab} — \textquotedblleft{}because of\textquotedblright{}. It is \emph{bi-}, \textquotedblleft{}by\textquotedblright{}, with \emph{sabab}, \textquotedblleft{}a reason\textquotedblright{}, and a noun follows it:

\begin{quote}
\textbf{\ar{تأخرت} \ar{بسبب} \ar{المطر}.} — \emph{taʾakhkhartu bi-sabab al-maṭar} — I was late because of the rain.
\end{quote}

Where the reason is a whole sentence, with its own verb, it takes \emph{li-anna}:

\begin{quote}
\textbf{\ar{تأخرت} \ar{لأن} \ar{الطريق} \ar{كان} \ar{مزدحما}.} — \emph{taʾakhkhartu li-anna aṭ-ṭarīqa kāna muzdaḥiman} — I was late because the road was busy.
\end{quote}

So \emph{bi-sabab al-maṭar}, a thing; \emph{li-anna aṭ-ṭarīqa kāna muzdaḥiman}, a sentence.
\end{grammarlens}

\begin{grammarlens}[title={Grammar Lens: the reason is that}]
To give the reason as an answer of its own, say \textbf{\ar{السبب} \ar{هو} \ar{أن}} — \emph{as-sababu huwa anna} — \textquotedblleft{}the reason is that\textquotedblright{}. \emph{Huwa}, \textquotedblleft{}he, it\textquotedblright{}, stands where English has \textquotedblleft{}is\textquotedblright{}:

\begin{quote}
\textbf{\ar{السبب} \ar{هو} \ar{أن} \ar{القطار} \ar{تأخر}.} — \emph{as-sababu huwa anna al-qiṭāra taʾakhkhara} — The reason is that the train was late.
\end{quote}

\emph{Natīja}, \textquotedblleft{}a result\textquotedblright{}, is its partner: \emph{as-sabab} looks back, \emph{an-natīja} looks forward.
\end{grammarlens}

\subsection*{Your turn}
Say these aloud:

\begin{itemize}
\raggedright
  \item \textquotedblleft{}I was late because of the rain\textquotedblright{} — \emph{taʾakhkhartu bi-sabab al-maṭar}
  \item \textquotedblleft{}because the road was busy\textquotedblright{} — \emph{li-anna aṭ-ṭarīqa kāna muzdaḥiman}
  \item \textquotedblleft{}The reason is that the train was late\textquotedblright{} — \emph{as-sababu huwa anna al-qiṭāra taʾakhkhara}
  \item a question and its answer — \emph{limādhā taʾakhkharta?} … \emph{bi-sabab al-qiṭār}
\end{itemize}

\begin{tcolorbox}[breakable,colback=teal!4,colframe=teal!35!black,title={Before you move on}]
What follows \emph{bi-sabab}? (\textbf{A noun}: \emph{bi-sabab al-maṭar}.) What follows \emph{li-anna}? (\textbf{A sentence.}) \textquotedblleft{}The reason is that\textquotedblright{}? (\emph{\textbf{As-sababu huwa anna.}}) Which word sits inside \emph{bi-sabab}? (\emph{\textbf{Sabab}}, a reason.)
\end{tcolorbox}
\section[\texorpdfstring{\ar{سوف} (sawfa)}{سوف (sawfa)}]{\ar{سوف} (sawfa) — will, for a plan}

\label{lesson:AR-C286-sawfa}



\begin{quote}
Say \textquotedblleft{}because of the rain\textquotedblright{}. (\emph{Bi-sabab al-maṭar}.) Then \textquotedblleft{}the reason is that\textquotedblright{}. (\emph{As-sababu huwa anna}.)

You have heard \emph{sa-usāfiru al-usbūʿa al-qādim}, \textquotedblleft{}I will travel next week\textquotedblright{}. Reasons are often reasons for a plan, and a plan needs a future.
\end{quote}

\begin{grammarlens}[title={Grammar Lens: sa- and sawfa}]
Arabic builds the future on the present. Put \textbf{\ar{سـ}} — \emph{sa-} — onto its front, with no space, or put \textbf{\ar{سوف}} — \emph{sawfa} — before it as a word of its own:

\begin{quote}
\textbf{\ar{سأسافر} \ar{غدا}.} — \emph{sa-usāfiru ghadan} — I will travel tomorrow.
\end{quote}

\begin{quote}
\textbf{\ar{سوف} \ar{نسافر} \ar{في} \ar{العطلة} \ar{القادمة}.} — \emph{sawfa nusāfiru fī al-ʿuṭla al-qādima} — We will travel in the next holiday.
\end{quote}

The two say the same \textquotedblleft{}will\textquotedblright{}. \emph{Sawfa} is the fuller one, and it is often heard for something further off or said with weight.
\end{grammarlens}

\begin{grammarlens}[title={Grammar Lens: lan, will not}]
For \textquotedblleft{}will not\textquotedblright{}, the \emph{sa-} goes and \textbf{\ar{لن}} — \emph{lan} — takes its place. After \emph{lan} the verb is in the subjunctive, as after \emph{an}:

\begin{quote}
\textbf{\ar{لن} \ar{أسافر} \ar{في} \ar{العطلة}.} — \emph{lan usāfira fī al-ʿuṭla} — I won't travel in the holiday.
\end{quote}

\emph{Lan adhhaba}, I won't go; \emph{lan nadhhaba}, we won't go.
\end{grammarlens}

\subsection*{Your turn}
Say these aloud:

\begin{itemize}
\raggedright
  \item \textquotedblleft{}I will travel tomorrow\textquotedblright{} — \emph{sa-usāfiru ghadan}
  \item \textquotedblleft{}We will travel in the next holiday\textquotedblright{} — \emph{sawfa nusāfiru fī al-ʿuṭla al-qādima}
  \item \textquotedblleft{}I won't go\textquotedblright{} — \emph{lan adhhaba}
  \item a plan with its reason — \emph{sa-usāfiru li-anna ʿindī waqtan}
  \item the opposite — \emph{lan usāfira li-anna ʿindī ʿamalan}
\end{itemize}

\begin{tcolorbox}[breakable,colback=teal!4,colframe=teal!35!black,title={Before you move on}]
The short \textquotedblleft{}will\textquotedblright{}? (\textbf{\ar{سـ}}, \emph{sa-}, joined to the verb.) The long one? (\emph{\textbf{Sawfa.}}) \textquotedblleft{}Will not\textquotedblright{}? (\emph{\textbf{Lan}}, with the subjunctive: \emph{lan adhhaba}.)
\end{tcolorbox}
\section[(review)]{Review — a plan, some advice, and why}

\label{lesson:AR-R286-advice-and-reasons}



\begin{quote}
Say \textquotedblleft{}I have a question\textquotedblright{}, then \textquotedblleft{}I think that it is easy\textquotedblright{}. (\emph{ʿIndī suʾāl}, \emph{aʿtaqidu annahu sahl}.) Then \textquotedblleft{}so\textquotedblright{}, the one-letter way. (\emph{Fa-}.)
\end{quote}

\begin{grammarlens}[title={What you've built}]
Four lessons, and one vowel at the end of a verb. After \emph{an}, after \emph{yajibu an} and \emph{yanbaghī an}, and after \emph{lan}, the present ends in \emph{-a}: the shape for a thing wanted, needed or refused, and not yet a fact. \emph{Sa-} and \emph{sawfa} keep the plain present. \emph{Bi-sabab} gives a reason that is a thing, \emph{li-anna} one that is a sentence.

A need and its reason in one breath:

\begin{quote}
\textbf{\ar{يجب} \ar{أن} \ar{أزور} \ar{الطبيب} \ar{غدا،} \ar{لأن} \ar{عندي} \ar{سؤالا}.} — \emph{yajibu an azūra aṭ-ṭabība ghadan, li-anna ʿindī suʾālan} — I must visit the doctor tomorrow, because I have a question.
\end{quote}
\end{grammarlens}

\subsection*{Your turn}
Say these aloud:

\begin{itemize}
\raggedright
  \item \textquotedblleft{}I want to travel\textquotedblright{}, then \textquotedblleft{}I won't travel\textquotedblright{} — \emph{urīdu an usāfira}, \emph{lan usāfira}
  \item \textquotedblleft{}You ought to rest\textquotedblright{} — \emph{yanbaghī an tastarīḥa}
  \item \textquotedblleft{}The lawyer must come now\textquotedblright{} — \emph{yajibu an yajīʾa al-muḥāmī al-ān}
  \item \textquotedblleft{}The reason is that the calculator is not here\textquotedblright{} — \emph{as-sababu huwa anna al-ḥāsibata laysat hunā}
  \item a plan of your own with \emph{sa-}, and its reason with \emph{li-anna} or \emph{bi-sabab}
\end{itemize}

\begin{tcolorbox}[breakable,colback=teal!4,colframe=teal!35!black,title={Before you move on}]
Which four words hand their verb the \emph{-a}? (\emph{\textbf{An}}, \emph{\textbf{yajibu an}}, \emph{\textbf{yanbaghī an}}, \emph{\textbf{lan}}.) \textquotedblleft{}I will go\textquotedblright{}? (\emph{\textbf{Sa-adhhabu.}}) \textquotedblleft{}Because of the rain\textquotedblright{}? (\emph{\textbf{Bi-sabab al-maṭar.}})
\end{tcolorbox}
